\documentclass[12pt]{article}
\usepackage[utf8]{inputenc}
\usepackage[T1]{fontenc}
\usepackage{hyperref}
\title{Khwarezmian Language Resources for K-12 Curriculum Development}
\author{Compiled by Claude Sonnet 5.5 (claude-sonnet-5-5) for Andr\'as Kornai}
\date{}
\begin{document}
\maketitle

\section{Introduction}
Khwarezmian (Chorasmian; ISO 639-3 xco; Glottocode khwa1238) is an extinct Eastern Iranian language closely related to Sogdian. It was spoken in Khwarezm, on the lower Amu Darya south of the Aral Sea, and was in use until at least the 13th century. Sources include astronomical terms in al-Biruni, Zamakhshari's Arabic--Persian--Khwarezmian dictionary and legal texts such as the \emph{Qunyat al-Munya}. I found no school materials.

\section{Primary Educational Resources}
I found no primers, textbooks or school programmes for Khwarezmian.

\section{Linguistic and Reference Materials}
\begin{itemize}
\item Wikipedia, ``Khwarezmian language'': overview of the language and its sources. \url{https://en.wikipedia.org/wiki/Khwarezmian_language}
\item Chorasmian Online (UC Berkeley, MELC): project site for the Chorasmian Dictionary Project. \url{https://chorasmianonline.melc.berkeley.edu}
\item Chorasmian Online, online dictionary. \url{https://chorasmianonline.melc.berkeley.edu/dictionary}
\item The Chorasmian Dictionary Project: Materials \& Goals. \url{https://chorasmianonline.melc.berkeley.edu/2024/01/07/chorasmian-dictionary-project-materials-goals}
\item Glottolog record for Henning's work, published by MacKenzie in 1971 as \emph{A Fragment of a Khwarezmian Dictionary}. Henning's entries run only from \textit{ʾ-} to \textit{ʾkw-}. \url{https://glottolog.org/resource/reference/id/58258}
\end{itemize}

\section{Religious and Community Resources}
I found no community organisations. The surviving texts include legal material (\emph{Qunyat al-Munya}). I did not find a URL for it in search results.

\section{Digital and Audio Resources}
The Chorasmian Online dictionary is the only digital resource I found. I found no audio.

\section{Curriculum Development Context}
The language is extinct and its record is fragmentary. The dictionary was never completed, and I did not find a dedicated grammar in the search results. A K-12 unit could treat Khwarezmian as a window on Central Asian history and on the Iranian language family, using the Berkeley project materials. Full language-teaching materials do not exist.
\end{document}